\documentclass[10pt,aspectratio=169]{beamer}

% --- PACOTES E CONFIGURAÇÕES DE IDIOMA E FONTE ---
\usepackage[utf8]{inputenc}
\usepackage[portuguese]{babel}
\usepackage[T1]{fontenc}
\usepackage{lmodern}
\usepackage{amsmath,amssymb}
\usepackage{graphicx}
\usepackage{booktabs}
\usepackage{colortbl}

% --- TEMA MADRID + WHALE ---
\usetheme{Madrid}
\usecolortheme{whale}

% Customização de cores adicionais
\definecolor{darkblue}{rgb}{0.1, 0.2, 0.4}
\definecolor{remlgreen}{rgb}{0.15, 0.55, 0.25}

\setbeamertemplate{navigation symbols}{} % Remove ícones de navegação poluídos

% --- INFORMAÇÕES DO DOCUMENTO ---
\title[Estatística Espacial]{Estimação de Parâmetros Espaciais: REML vs Variograma}
\subtitle{Avaliação de Desempenho e Eficiência via Simulação Monte Carlo}
\author{Azzure Alves do Carmo}
\institute[UFES]{Baseado em \textit{Spatial Statistics for Data Science: Theory and Practice with R}}
\date{\today}

\begin{document}

% ------------------------------------------------------------------------------
% SLIDE 1: TÍTULO
% ------------------------------------------------------------------------------
\begin{frame}
  \titlepage
\end{frame}

% ------------------------------------------------------------------------------
% SLIDE 2: CONTEXTO E OBJETIVO ESTRATÉGICO
% ------------------------------------------------------------------------------
\begin{frame}{Contexto e Objetivo}
  \begin{block}{Motivação}
    \begin{itemize}
      \item A \textbf{Krigagem Ordinária} fornece predições espaciais ótimas de variância mínima (Melhor Preditor Linear Não Viesado - BLUP) cujos pesos dependem da função de covariância espacial.
      
      \item Tradicionalmente, os parâmetros (efeito pepita $\tau^2$, patamar $\sigma^2$ e alcance $\phi$) são estimados por variograma, usando Mínimos Quadrados Ponderados - WLS e Mínimos Quadrados Ordinários - OLS, ou por verossimilhança (Máxima Verossimilhança Restrita - REML).
    \end{itemize}
  \end{block}

  \vspace{0.2cm}

  \begin{block}{Objetivo Principal}
    Comparar a acurácia (viés) e a precisão (Erro Quadrático Médio - RMSE) da Máxima Verossimilhança Restrita (REML) contra os métodos baseados no variograma (WLS e OLS) sob diferentes tamanhos amostrais ($n \in \{50, 100, 200\}$).
  \end{block}
\end{frame}

% ------------------------------------------------------------------------------
% SLIDE 3: FORMULAÇÃO TEÓRICA DO REML
% ------------------------------------------------------------------------------
\begin{frame}{Máxima Verossimilhança Restrita (REML)}
  \begin{block}{Modelo Gaussiano Espacial}
    Para um processo espacial $Y \sim \mathcal{N}(X\beta, \Sigma(\theta))$, onde $\Sigma(\theta) = \sigma^2 R(\phi) + \tau^2 I_n$:
  \end{block}

  \vspace{0.1cm}

  \begin{block}{Log-Verossimilhança Restrita}
    O método REML separa a estimação dos parâmetros de covariância $\theta = (\tau^2, \sigma^2, \phi)^T$ da estimação dos efeitos fixos $\beta$, eliminando o viés de pequena amostra:
    \begin{equation*}
      l_{\text{REML}}(\theta) = -\frac{1}{2} \log |\Sigma(\theta)| - \frac{1}{2} \log |X^T \Sigma(\theta)^{-1} X| - \frac{1}{2} (y - X\hat{\beta})^T \Sigma(\theta)^{-1} (y - X\hat{\beta})
    \end{equation*}
    onde $\hat{\beta} = (X^T \Sigma(\theta)^{-1} X)^{-1} X^T \Sigma(\theta)^{-1} y$ é o estimador de Mínimos Quadrados Generalizados (GLS).
  \end{block}

  \begin{block}{Vantagem Teórica}
    Ao contrário da Máxima Verossimilhança direta (ML), o REML corrige os graus de liberdade perdidos na estimação da média $\beta$.
  \end{block}
\end{frame}

% ------------------------------------------------------------------------------
% SLIDE 4: DESENHO EXPERIMENTAL
% ------------------------------------------------------------------------------
\begin{frame}{Metodologia da Simulação Monte Carlo}
  \begin{columns}[T]
    \begin{column}{0.48\textwidth}
      \begin{block}{Parâmetros Verdadeiros}
        \begin{itemize}
          \item \textbf{Efeito Pepita ($\tau^2$):} $0.20$ (Ruído em microescala)
          \item \textbf{Patamar ($\sigma^2$):} $1.00$ (Variância espacial)
          \item \textbf{Alcance ($\phi$):} $20.0$ (Decaimento espacial)
          \item \textbf{Modelo:} Exponencial $\gamma(h) = \tau^2 + \sigma^2(1 - e^{-h/\phi})$
        \end{itemize}
      \end{block}
    \end{column}

    \begin{column}{0.48\textwidth}
      \begin{block}{Configuração da Simulação}
        \begin{itemize}
          \item \textbf{Réplicas Monte Carlo:} $N = 100$ por cenário
          \item \textbf{Tamanhos Amostrais ($n$):} $50, 100, 200$ pontos
          \item \textbf{Domínio Espacial:} $[0, 100] \times [0, 100]$
          \item \textbf{Métodos Comparados:}
            \begin{enumerate}
              \item \textbf{REML}: Verossimilhança exata no subespaço
              \item \textbf{WLS}: Pesos de Cressie $w_k = \frac{N(h_k)}{\gamma(h_k)^2}$
              \item \textbf{OLS}: Pesos unitários $w_k = 1$
            \end{enumerate}
        \end{itemize}
      \end{block}
    \end{column}
  \end{columns}
\end{frame}

% ------------------------------------------------------------------------------
% SLIDE 5: EXEMPLO DE VARIOGRAMA E AJUSTES
% ------------------------------------------------------------------------------
\begin{frame}{Exemplo de Variograma Empírico e Ajustes}
  \begin{center}
    \includegraphics[height=0.73\textheight,keepaspectratio]{plots/exemplo_variograma_ajustado.png}
  \end{center}
\end{frame}

% ------------------------------------------------------------------------------
% SLIDE 6: DISTRIBUIÇÃO DOS PARÂMETROS (REML vs WLS vs OLS)
% ------------------------------------------------------------------------------
\begin{frame}{Distribuição dos Parâmetros: REML vs WLS vs OLS}
  \begin{center}
    \includegraphics[height=0.73\textheight,keepaspectratio]{plots/boxplot_reml_vs_wls.png}
  \end{center}
\end{frame}

% ------------------------------------------------------------------------------
% SLIDE 7: EVOLUÇÃO DO RMSE (REML vs WLS vs OLS)
% ------------------------------------------------------------------------------
\begin{frame}{Comparação de Eficiência (RMSE por Amostra)}
  \begin{center}
    \includegraphics[height=0.73\textheight,keepaspectratio]{plots/rmse_reml_vs_wls.png}
  \end{center}
\end{frame}

% ------------------------------------------------------------------------------
% SLIDE 8: TABELA RESUMO DE RESULTADOS
% ------------------------------------------------------------------------------
\begin{frame}{Resultados Numéricos Comparativos ($N = 100$)}
  \scriptsize
  \centering
  \setlength{\tabcolsep}{4.5pt}
  \begin{table}[h]
    \centering
    \caption{Desempenho dos Estimadores: Média Estimada e RMSE por Método}
    \begin{tabular}{ccccccccc}
      \toprule
      & & & \multicolumn{2}{c}{\textbf{REML}} & \multicolumn{2}{c}{\textbf{WLS}} & \multicolumn{2}{c}{\textbf{OLS}} \\
      \cmidrule(lr){4-5} \cmidrule(lr){6-7} \cmidrule(lr){8-9}
      \textbf{$n$} & \textbf{Parâmetro} & \textbf{Real} & \textbf{Média} & \textbf{RMSE} & \textbf{Média} & \textbf{RMSE} & \textbf{Média} & \textbf{RMSE} \\
      \midrule
      $n = 50$ & Efeito Pepita ($\tau^2$) & 0.20 & \textbf{0.186} & \textbf{0.177} & 0.287 & 0.257 & 0.218 & 0.198 \\
      $n = 50$ & Patamar ($\sigma^2$)     & 1.00 & \textbf{1.191} & \textbf{0.931} & 1.468 & 1.496 & 1.430 & 1.400 \\
      $n = 50$ & Alcance ($\phi$)         & 20.0 & \textbf{33.00} & \textbf{43.10} & 50.93 & 62.56 & 43.76 & 54.62 \\
      \midrule
      $n = 100$ & Efeito Pepita ($\tau^2$) & 0.20 & \textbf{0.187} & \textbf{0.107} & 0.237 & 0.145 & 0.208 & 0.138 \\
      $n = 100$ & Patamar ($\sigma^2$)     & 1.00 & \textbf{1.145} & \textbf{0.578} & 1.395 & 1.116 & 1.376 & 1.133 \\
      $n = 100$ & Alcance ($\phi$)         & 20.0 & \textbf{26.96} & \textbf{26.80} & 42.75 & 52.42 & 38.13 & 46.74 \\
      \midrule
      $n = 200$ & Efeito Pepita ($\tau^2$) & 0.20 & \textbf{0.192} & \textbf{0.068} & 0.245 & 0.119 & 0.231 & 0.109 \\
      $n = 200$ & Patamar ($\sigma^2$)     & 1.00 & \textbf{1.002} & \textbf{0.357} & 1.097 & 0.742 & 1.068 & 0.625 \\
      $n = 200$ & Alcance ($\phi$)         & 20.0 & \textbf{20.30} & \textbf{10.48} & 31.16 & 36.52 & 28.32 & 30.68 \\
      \bottomrule
    \end{tabular}
  \end{table}
\end{frame}

% ------------------------------------------------------------------------------
% SLIDE 9: RECOMENDAÇÕES GERENCIAIS
% ------------------------------------------------------------------------------
\begin{frame}{Conclusões}
  \begin{block}{Principais Descobertas}
    \begin{itemize}
      \item O REML foi consideravelmente melhor do que WLS e OLS em acurácia e precisão (menor RMSE) em todos os cenários amostrais testados.
      \item Portanto, para amostras pequenas a médias ($n \le 200$), priorizar a estimação por REML devido à eficiência estatística superior. Para aplicações de grande escala ($n > 1000$), pode ser necessário utilizar WLS.
    \end{itemize}
  \end{block}
\end{frame}

\end{document}
